%% 
%% Second Paper — Elsevier CAS Double Column
%% Delay-Integrated Cascading Failure Analysis in CPPS
%% 
\documentclass[a4paper,fleqn]{cas-dc}

\usepackage[numbers]{natbib}
\usepackage{amsmath,amssymb}
\usepackage{booktabs}
\usepackage{multirow}
\usepackage{graphicx}
\usepackage{xcolor}
%% \usepackage{algorithm}
%% \usepackage{algorithmic}

%% Custom commands
\newcommand{\Rone}{R_1}
\newcommand{\Rthree}{R_3}
\newcommand{\nodelay}{\texttt{no\_delay}}
\newcommand{\tauCC}{\tau_{\text{CC}}}
\newcommand{\tauNCC}{\tau_{\text{nonCC}}}

\begin{document}
\let\WriteBookmarks\relax
\def\floatpagepagefraction{1}
\def\textpagefraction{.001}

%% ====================================================================
%%  Front matter
%% ====================================================================
\shorttitle{Delay-Integrated Cascading Failure in CPPS}
\shortauthors{Z. Chen et~al.}

\title[mode=title]{Delay-Integrated Cascading Failure Analysis in
  Cyber-Physical Power Systems: A Closed-Loop Feedback Perspective}

\author[1]{Zhange Chen}
\author[1]{Haicheng Tu}\cormark[1]
\ead{tuhc@hdu.edu.cn}
\author[1]{Yongxiang Xia}
\author[2]{Xuetao Yang}

\cortext[1]{Corresponding author.}

\affiliation[1]{organization={School of Communication Engineering, Hangzhou Dianzi University},
  city={Hangzhou}, country={China}}
\affiliation[2]{organization={Power System Department, China Electric Power Research Institute},
  city={Beijing}, country={China}}

\credit{Z.~Chen: Methodology, Software, Writing -- Original Draft.
  H.~Tu: Conceptualization, Supervision, Writing -- Review \& Editing.
  Y.~Xia: Formal Analysis, Validation.
  X.~Yang: Resources, Funding Acquisition.}

%% ---- Abstract --------------------------------------------------------
\begin{abstract}
Cascading failure models for cyber-physical power systems (CPPS) routinely
treat communication delay as a post-hoc observation rather than a causal
driver of further damage.  This architectural choice eliminates the
mechanism through which delay inflicts its greatest harm: a closed
feedback loop in which degraded communication reduces generator
controllability, the resulting power imbalance triggers additional branch
overloads, the consequent structural failures further damage the
communication topology, and the damaged topology amplifies delay for the
next cascade round.

In this paper we embed communication delay directly into the cascading
failure process.  At each cascade round, a delay-dependent control
efficiency factor is computed for every surviving generator and applied
before the DC power flow that determines overloads.  End-to-end latency
is decomposed into uplink measurement delay and downlink execution delay
through a role-aware model that distinguishes control centers from
ordinary relay nodes.  Interdependent coupling between the IEEE~39-bus
power grid and a Barab\'{a}si--Albert scale-free communication network is
established via a betweenness-assortative strategy, with control centers
placed using a composite structural importance score.

Monte Carlo experiments spanning 1\,000 independent coupling realizations,
eleven tolerance levels, and five delay-severity scenarios yield three
principal findings: (i)~the delay penalty concentrates asymmetrically at
moderate-to-high tolerance margins---precisely where delay-free analysis
predicts safety but delay-burdened systems collapse; (ii)~delay amplifies
the stochastic variability of cascade outcomes, transforming narrow
distributions under ideal communication into heavy-tailed distributions
under severe delay; and (iii)~generator dispatch deviation remains
stratified across delay scenarios even when aggregate load retention
appears adequate, exposing hidden functional degradation invisible to
load-based metrics alone.  These results demonstrate that delay-free
robustness models are not merely imprecise but systematically misleading,
providing the greatest reassurance in operating regimes that harbor the
greatest vulnerability.
\end{abstract}

%% ---- Highlights ------------------------------------------------------
\begin{highlights}
\item A closed-loop cascade framework embeds communication delay as an
  endogenous causal variable inside each cascade round.
\item Delay damage is asymmetrically concentrated at moderate-to-high
  tolerance margins where delay-free models predict safety.
\item Delay amplifies stochastic variability, creating heavy-tailed
  collapse distributions absent under ideal communication.
\item Generator dispatch deviation metric $R_3$ reveals hidden
  functional degradation invisible to aggregate load-retention metrics.
\item Round-level heatmap analysis identifies Round~2 as the critical
  intervention window for delay mitigation.
\end{highlights}

%% ---- Keywords --------------------------------------------------------
\begin{keywords}
Cyber-physical power system \sep Cascading failure \sep
Communication delay \sep Closed-loop feedback \sep
Robustness assessment \sep Interdependent networks
\end{keywords}

\maketitle

%% ====================================================================
%%  SECTION 1 — INTRODUCTION
%% ====================================================================
\section{Introduction}\label{sec:intro}

The deep integration of physical infrastructure with communication
networks has transformed traditional power grids into cyber-physical
power systems (CPPS), in which a physical power layer and a cyber
communication layer operate as a tightly coupled whole
\citep{Rinaldi2001,Ouyang2014}.  Under normal conditions, this coupling
improves dispatch flexibility and situational awareness.  Under stress,
however, it creates a bidirectional failure channel: a disturbance
originating in either layer can propagate to the other, triggering
cascading failures that neither layer would experience in isolation
\citep{Buldyrev2010,Gao2012}.

A rich body of literature has investigated cascading failures in coupled
networks through topological models, flow-redistribution rules, and
attack-vulnerability analyses
\citep{Motter2002,Dobson2007,Brummitt2012}.  These studies have
established how structural interdependence amplifies failure propagation.
Yet most of them treat the communication layer as a binary entity---a
link is either functional or failed---without modeling the \emph{quality}
of communication service that surviving links provide.  In particular,
communication delay, the time elapsed between issuing a measurement or
command and its arrival at the destination, is largely absent from
cascading-failure analysis.  When delay is mentioned, it is typically
reported as a post-hoc observation (``delay increased after failures'')
rather than modeled as a causal variable that directly degrades generator
control performance during the cascade itself
\citep{Milano2012,Huang2015}.

This omission matters more than it might first appear.  In a closed-loop
control architecture, stale measurements cause the control center to
make decisions based on outdated system states, while lagged command
execution means that corrective actions arrive after the system has
drifted further.  Both effects reduce the effective generation output
relative to its intended reference.  When this reduction is embedded
directly into the cascading failure process---before the power-flow
redistribution that determines which branches overload---a
self-reinforcing feedback loop emerges:
\begin{quote}
\emph{Delay reduces generation $\;\to\;$ altered power flows cause
additional overloads $\;\to\;$ overloaded branches trip $\;\to\;$
structural failures degrade the cyber topology $\;\to\;$ degraded
topology increases delay $\;\to\;$ further output reduction
$\;\to\;\cdots$}
\end{quote}
This feedback loop transforms communication latency from a passive
symptom of system distress into an active driver of further damage.
Crucially, it is invisible to models that apply delay as a post-hoc
correction, in which all delay scenarios share an identical structural
cascade and delay merely rescales final performance metrics
\citep{LaiCPPS2019}.

Our prior work \citep{ChenTu2026} established a CPPS cascading failure
framework with distributed control centers, a Functional
Viability-Constrained Island Detection (FVC-ID) mechanism, and a
Cohesion-Propagation Importance Metric (CPIM) for optimal control-center
placement.  That study did not, however, model communication delay's
influence on generator control during the cascade.  The present paper
fills this gap.

The contributions of this work are as follows:

\begin{enumerate}
\item \textbf{A delay-integrated cascading failure framework with
  closed-loop feedback.}  We embed delay computation directly into each
  cascade round, scaling generator output before DC power flow
  redistribution.  Each delay scenario therefore produces a genuinely
  different cascade trajectory---different overload patterns, different
  failure sequences, different terminal states---rather than a rescaled
  version of a common trajectory.

\item \textbf{A bidirectional, role-aware delay model.}  Uplink
  (measurement) and downlink (command) delays are separately computed
  from link-level serial and propagation components, with
  role-dependent service delays for control centers versus ordinary
  relay nodes.  This captures the asymmetric processing behavior
  inherent in SCADA-type control architectures.

\item \textbf{Mechanism decomposition into two degradation channels.}
  We decompose delay-induced robustness loss into continuous efficiency
  erosion (path lengthening reduces~$\eta$) and discontinuous
  reachability collapse (cyber-network fragmentation severs all CC
  access, setting~$\eta=0$).

\item \textbf{Identification of the asymmetric delay-penalty regime.}
  Through a delay-penalty heatmap over the (tolerance, cascade-round)
  plane, we demonstrate that delay damage is largest at moderate-to-high
  tolerance margins---where the delay-free system retains most load but
  the delay-burdened system collapses---revealing that delay-free models
  create a false sense of security precisely where operators believe the
  system is adequately protected.
\end{enumerate}

The remainder of this paper is organized as follows.
Section~\ref{sec:model} presents the CPPS model, including the
communication delay formulation and interdependent coupling.
Section~\ref{sec:cascade} describes the delay-integrated cascading
failure dynamics.  Section~\ref{sec:experiment} reports experimental
results.  Section~\ref{sec:conclusion} concludes the paper.


%% ====================================================================
%%  SECTION 2 — CPPS MODEL
%% ====================================================================
\section{Cyber-Physical Power System Model}\label{sec:model}

The system under study consists of two interdependent network layers---a
physical power grid and a cyber communication network---whose coupling
transforms isolated disturbances into cross-domain cascading failures.

%% ---- 2.1 Power Grid Layer -------------------------------------------
\subsection{Power Grid Layer}\label{sec:power}

The physical layer is represented by the IEEE~39-bus (New England) test
system, a well-established benchmark comprising 39~buses,
46~transmission branches, and 10~generators.  Bus connectivity is
encoded in an adjacency matrix $A_P$ extracted from the branch table,
and the initial operating state---bus loads, generator outputs, and
branch power flows---is obtained by solving a DC power flow on the full
network \citep{Zimmerman2011matpower}.

Under the DC power flow assumption, the active power injections satisfy
\begin{equation}\label{eq:dcpf}
  \mathbf{P} = \mathbf{B}\boldsymbol{\Theta},
\end{equation}
where $\mathbf{P}=[p_1,p_2,\dots,p_n]^{T}$ is the nodal injection
vector ($p_i>0$ for generators, $p_i<0$ for loads),
$\boldsymbol{\Theta}$ is the bus voltage-angle vector, and $\mathbf{B}$
is the bus susceptance matrix with elements
\begin{equation}
  B_{ij}=\begin{cases}
    -b_{ij}, & i\ne j,\\[4pt]
    \displaystyle\sum_{k\ne i}b_{ik}, & i=j,
  \end{cases}
\end{equation}
where $b_{ij}=1/x_{ij}$ is the reciprocal of the reactance of the
line connecting buses $i$ and~$j$.  Each branch~$e$ carries an initial
power flow $P_e^{(0)}$ that serves as the reference for the capacity
model in Section~\ref{sec:capacity}.

%% ---- 2.2 Communication Network Layer --------------------------------
\subsection{Communication Network Layer}\label{sec:cyber}

\subsubsection{Scale-Free Topology}

The cyber layer is constructed as a Barab\'{a}si--Albert (BA) scale-free
network to capture the hub-dominated, heterogeneous degree distribution
commonly observed in communication infrastructures
\citep{BarabasiAlbert1999}.  The network contains
$|V_C|=|V_P|+n_{\text{cc}}$ nodes, where
$n_{\text{cc}}=\lceil 0.2\times|V_P|\rceil$ additional nodes serve as
designated control centers (CCs).  For the IEEE~39-bus system this
yields $|V_C|=39+8=47$ nodes, of which 8 are CCs.

The BA growth process starts from a fully connected seed of $m_0=4$
nodes.  At each step a new node is added and connected to $m_e=2$
existing nodes selected with probability proportional to their current
degree (preferential attachment), producing a power-law degree
distribution characteristic of real communication networks.

\subsubsection{Control Center Placement via MSIS}\label{sec:msis}

Rather than selecting CCs by a single centrality measure, we employ a
Multi-attribute Structural Importance Score (MSIS) that fuses
complementary topological features:
\begin{equation}\label{eq:msis}
  w_i = \frac{\widetilde{B}(i)+K_s(i)}{C_l(i)} + D(i)\cdot C(i),
\end{equation}
where $\widetilde{B}(i)=B(i)+1$ is the shifted betweenness centrality,
$K_s(i)$ is the $k$-shell index, $C_l(i)$ is closeness centrality,
$D(i)$ is degree, and $C(i)$ is the local clustering coefficient.  The
top-$n_{\text{cc}}$ nodes ranked by $w_i$ are designated as control
centers.

This composite score balances global bridging importance (betweenness,
closeness) with local structural embeddedness (degree, clustering,
$k$-shell).  Our prior work proposed the Cohesion-Propagation Importance
Metric (CPIM) for the same purpose \citep{ChenTu2026}; MSIS adopts a
similar multi-attribute philosophy while using a different functional
form suited to the delay-focused analysis of this paper.

\subsubsection{Communication Delay Model}\label{sec:delay_model}

Communication delay is modeled at the individual link level and
aggregated over multi-hop paths, with explicit directional asymmetry
between uplink (measurement) and downlink (command) traffic.

\paragraph{Link-level delay.}
For a single surviving link $(u,v)$, the per-hop delay is
\begin{equation}\label{eq:link_delay}
  T_{\text{link}}(u,v) = \frac{S}{R} + \frac{d}{V},
\end{equation}
where $S$ is the packet size (bits), $R$ is the link data rate~(bps),
$d$ is the effective link distance~(km), and $V$ is the signal
propagation speed~(km/s).  The first term captures serialization delay;
the second captures propagation delay.  Uplink and downlink traffic use
different packet sizes---measurement packets ($S_{\text{up}}=8\,192$~bits)
are larger than command packets ($S_{\text{down}}=2\,048$~bits)---introducing
directional asymmetry even over the same physical link.

\paragraph{Path-level delay.}
For a multi-hop path $\pi=(u_0,u_1,\dots,u_k)$, the total path delay
aggregates per-hop link delays, endpoint service delays, and
intermediate forwarding delays.

\textit{Downlink} (CC $\to$ non-CC):
\begin{equation}\label{eq:Tdown}
  T_{\text{down}}(\pi) = \sum_{i=0}^{k-1}T_{\text{link}}^{\text{down}}(u_i,u_{i+1})
  + \tauCC^{\text{tx}} + \tauNCC^{\text{rx}}
  + \sum_{j=1}^{k-1}\tau_{\text{fwd}}(u_j).
\end{equation}

\textit{Uplink} (non-CC $\to$ CC):
\begin{equation}\label{eq:Tup}
  T_{\text{up}}(\pi) = \sum_{i=0}^{k-1}T_{\text{link}}^{\text{up}}(u_i,u_{i+1})
  + \tauNCC^{\text{tx}} + \tauCC^{\text{rx}}
  + \sum_{j=1}^{k-1}\tau_{\text{fwd}}(u_j).
\end{equation}

The forwarding delay at each intermediate node is role-dependent:
\begin{equation}\label{eq:fwd}
  \tau_{\text{fwd}}(u) =
  \begin{cases}
    \tauCC^{\text{fwd}},   & \text{if } u \text{ is a CC},\\
    \tauNCC^{\text{fwd}},  & \text{otherwise}.
  \end{cases}
\end{equation}
Control centers are equipped with more capable processing hardware and
therefore exhibit lower per-packet service latency than ordinary relay
nodes.  The same physical path thus yields different uplink and downlink
delays because of directional traffic characteristics and role-dependent
processing, not because of asymmetric propagation physics.

%% ---- 2.3 Interdependent Coupling ------------------------------------
\subsection{Interdependent Coupling}\label{sec:coupling}

The two layers are coupled through a one-to-one mapping between power
buses and non-CC cyber nodes, encoded in a coupling matrix
$A_{pc}\in\{0,1\}^{|V_P|\times|V_C|}$.

\paragraph{Betweenness-assortative coupling.}
We construct $A_{pc}$ by sorting power buses in descending order of
their betweenness centrality in~$G_P$, independently sorting non-CC
cyber nodes in descending order of their betweenness centrality
in~$G_C$, and pairing them rank-by-rank.  This assortative strategy
ensures that structurally critical power buses are paired with
structurally critical communication nodes, reflecting the engineering
practice of provisioning important substations with high-quality
communication resources \citep{ChenTu2026}.

Because each Monte Carlo trial generates a fresh BA communication
network, a fresh CC selection, and a fresh coupling matrix, the reported
metrics reflect averaged behavior across diverse cyber-physical
configurations.

%% ---- 2.4 Capacity-Tolerance Model -----------------------------------
\subsection{Capacity-Tolerance Model}\label{sec:capacity}

Overload-driven cascading requires a model of component capacity.  We
adopt the standard tolerance-factor formulation:

\textit{Power branches:}
\begin{equation}\label{eq:cap_P}
  C_e^{P} = (1+\alpha)\cdot|P_e^{(0)}|,
\end{equation}
where $P_e^{(0)}$ is the initial branch flow from DC power flow.

\textit{Cyber nodes:}
\begin{equation}\label{eq:cap_Cn}
  C_i^{C} = (1+\alpha)\cdot B_i^{(0)},
\end{equation}
where $B_i^{(0)}$ is the initial betweenness centrality of cyber node~$i$.

\textit{Cyber edges:}
\begin{equation}\label{eq:cap_Ce}
  C_e^{CE} = (1+\alpha)\cdot B_e^{(0)},
\end{equation}
where $B_e^{(0)}$ is the initial edge betweenness of cyber link~$e$.

The tolerance factor $\alpha\in[0,1]$ governs excess capacity.  At
$\alpha=0$ components operate at full capacity with no margin; at
$\alpha=1$ each component can absorb up to double its initial load.  By
sweeping $\alpha$ from 0 to~1 in increments of~0.1, we systematically
explore how the interplay between structural margin and communication
delay shapes cascade severity.


%% ====================================================================
%%  SECTION 3 — DELAY-INTEGRATED CASCADE DYNAMICS
%% ====================================================================
\section{Delay-Integrated Cascading Failure Dynamics}\label{sec:cascade}

The cascading failure model operates through a nested dual-loop
architecture in which communication delay participates as an endogenous
causal variable that actively shapes the cascade trajectory.

%% ---- 3.1 Attack Model -----------------------------------------------
\subsection{Attack Model}\label{sec:attack}

Each simulation begins with the targeted removal of a single non-CC
cyber node---the non-CC node with the highest betweenness centrality in
the initial communication network.  This betweenness-targeted attack
maximizes initial disruption to communication paths: the attacked node
is a high-betweenness relay whose removal simultaneously lengthens many
communication paths and may fragment the cyber network
\citep{Albert2000}.

%% ---- 3.2 Structural Propagation (Inner Loop) -----------------------
\subsection{Structural Propagation---Inner Loop}\label{sec:inner}

Given the current set of failed nodes and edges, the inner loop iterates
until no new structural failures emerge.  Each iteration executes four
steps:

\paragraph{Step~1: Cyber-layer island viability (FVC-ID).}
Connected components of the surviving cyber network are identified.  Any
component that does not contain at least one surviving control center is
declared non-viable---without CC access, nodes can neither report
measurements nor receive dispatch commands.  All nodes and edges within
non-viable components are marked as failed.  This implements the
Functional Viability-Constrained Island Detection (FVC-ID) mechanism
introduced in our prior work \citep{ChenTu2026}.

\paragraph{Step~2: Forward propagation (cyber $\to$ power).}
Each newly failed non-CC cyber node triggers the failure of its coupled
power bus with probability~$p$.

\paragraph{Step~3: Power-layer island viability.}
Connected components of the surviving power network are identified.  Any
component that contains neither a generator nor a load bus is declared
non-viable, and all its buses and branches are marked as failed.

\paragraph{Step~4: Reverse propagation (power $\to$ cyber).}
Each newly failed power bus triggers the failure of its coupled non-CC
cyber node with probability~$p$.

\paragraph{Convergence.}
Steps~1--4 repeat until the sets of failed cyber and power nodes remain
unchanged.  The bidirectional propagation with probability~$p$ is the
mechanism through which damage crosses the layer boundary, creating
cascading chains that alternate between layers until convergence
\citep{Parshani2010,Shao2011}.

%% ---- 3.3 Overload-Driven Cascade (Outer Loop) -----------------------
\subsection{Overload-Driven Cascade Expansion---Outer Loop}\label{sec:outer}

After the inner loop stabilizes, the surviving network undergoes
overload evaluation with delay injection.

\paragraph{Step~1: Cyber-layer overload check.}
Betweenness centrality is recomputed on the surviving cyber network.
Any non-CC node whose updated betweenness exceeds its capacity
$C_i^C=(1+\alpha)\cdot B_i^{(0)}$ is removed.  Similarly, any link
whose edge betweenness exceeds $C_e^{CE}=(1+\alpha)\cdot B_e^{(0)}$ is
removed.

\paragraph{Step~2: Delay-aware generation adjustment.}
For each surviving generator, the framework computes a delay-dependent
control efficiency factor~$\eta_i$ based on the current state of the
communication network and scales the generator's output accordingly.
This step---detailed in Section~\ref{sec:delay_dispatch}---is the
architectural innovation that closes the feedback loop.  The scaling
occurs \emph{before} the DC power flow, ensuring that delay's impact
propagates through the electrical network.

\paragraph{Step~3: DC power flow and branch overload.}
A DC power flow is solved using the delay-adjusted generator outputs.
Any branch whose post-redistribution flow exceeds
$C_e^P=(1+\alpha)\cdot|P_e^{(0)}|$ is tripped.

\paragraph{Step~4: Convergence test.}
If no new failures occur in Steps~1 and~3, the cascade terminates.
Otherwise, the newly failed components are incorporated and the process
returns to the inner loop (Section~\ref{sec:inner}).

\paragraph{The feedback mechanism.}
In a conventional cascade model, every generator produces its full
scheduled output regardless of communication quality.  In the present
framework, generators whose communication paths have lengthened produce
less power (reduced~$\eta_i$), and generators that have lost all CC
access produce no power ($\eta_i=0$).  This output reduction alters
branch flows: branches that would not overload under ideal control may
now overload because reduced generation forces more aggressive
redistribution onto surviving paths.  These additional overloads cause
further structural failures, which degrade the cyber topology, which
increases delay.  Each delay scenario therefore traces a unique cascade
trajectory---a property fundamentally impossible in post-hoc delay
models.

%% ---- 3.4 Delay-Aware Generator Dispatch -----------------------------
\subsection{Delay-Aware Generator Dispatch}\label{sec:delay_dispatch}

For each surviving generator~$i$ that remains online after structural
propagation:

\begin{enumerate}
\item \textbf{Identify the coupled cyber node.}  Generator~$i$'s power
  bus is mapped to its non-CC cyber node via~$A_{pc}$.

\item \textbf{Find the nearest reachable CC.}  On the surviving cyber
  network, the shortest delay-weighted path from the generator's cyber
  node to each surviving CC is computed.  The CC with the minimum-delay
  path is selected.  If no CC is reachable, the generator is classified
  as \emph{unreachable}.

\item \textbf{Compute delays.}  The total measurement delay and
  execution delay for generator~$i$ are
  \begin{align}
    \tau_{m,i} &= \tau_{\text{PB}\to\text{nonCC}}
                  + T_{\text{up}}(\pi_i), \label{eq:tau_m}\\
    \tau_{e,i} &= \tau_{\text{nonCC}\to\text{PB}}
                  + T_{\text{down}}(\pi_i), \label{eq:tau_e}
  \end{align}
  where $\tau_{\text{PB}\to\text{nonCC}}$ is the field-side measurement
  interface delay and $\tau_{\text{nonCC}\to\text{PB}}$ is the
  field-side command execution delay.

\item \textbf{Compute control efficiency.}  The delay-to-control
  mapping converts the two delay components into a composite efficiency
  factor:
  \begin{equation}\label{eq:eta}
    \eta_i = \min\!\bigl(1,\;\max\!\bigl(0,\;
      (1-k_m\tau_{m,i})(1-k_e\tau_{e,i})\bigr)\bigr),
  \end{equation}
  where $k_m$ and $k_e$ are delay sensitivity coefficients.

\item \textbf{Scale generator output.}  The actual output is
  \begin{equation}\label{eq:Pactual}
    P_i^{\text{actual}} = P_i^{\text{ref}}\cdot\eta_i.
  \end{equation}
\end{enumerate}

\paragraph{Unreachable generators.}
If no surviving CC is reachable, the generator receives no control
commands: $P_i^{\text{actual}}=0$, $\eta_i=0$.  This is the most severe
consequence of cyber-layer degradation---complete loss of
controllability---and it occurs before the DC power flow, directly
contributing to power imbalance and further overloads.

%% ---- 3.5 Delay Severity Scenarios -----------------------------------
\subsection{Delay Severity Scenarios}\label{sec:scenarios}

Five delay-severity scenarios are defined by coherently scaling a common
baseline parameter set (Table~\ref{tbl:scenarios}).  The scale factor is
applied multiplicatively to all delay-producing parameters: packet
sizes, link distances, all service delays, and field-side interface
delays.  The link data rate~$R$, propagation speed~$V$, and sensitivity
coefficients $k_m$, $k_e$ remain fixed.

\begin{table}[t]
\caption{Delay severity scenarios.}\label{tbl:scenarios}
\centering
\begin{tabular}{@{}lcc@{}}
\toprule
Scenario & Scale factor & Interpretation \\
\midrule
\texttt{no\_delay} & 0.0 & Zero communication delay \\
\texttt{light}     & 0.5 & Well-provisioned communication \\
\texttt{baseline}  & 1.0 & Nominal delay parameterization \\
\texttt{medium}    & 1.5 & Moderately degraded \\
\texttt{heavy}     & 2.0 & Severely degraded \\
\bottomrule
\end{tabular}
\end{table}

The \nodelay{} scenario (scale~0) zeros all delay terms; every reachable
generator retains $\eta_i=1$.  Any performance difference between
\nodelay{} and other scenarios is therefore causally attributable to
communication delay.

%% ---- 3.6 Resilience Metrics ----------------------------------------
\subsection{Resilience Metrics}\label{sec:metrics}

\paragraph{$\Rone$: Load retention ratio.}
Let $L_{\text{initial}}$ denote the total pre-attack load and
$L_{\text{surviving}}$ the total load at surviving buses:
\begin{equation}\label{eq:R1}
  \Rone = \frac{L_{\text{surviving}}}{L_{\text{initial}}}.
\end{equation}
Because delay is embedded within the cascade process, delay's influence
on~$\Rone$ manifests endogenously through divergent cascade trajectories.

\paragraph{$\Rthree$: Generator dispatch deviation.}
For the $N$ generators surviving at the terminal state with
$P_j^{\text{ref}}>0$:
\begin{equation}\label{eq:R3}
  \Rthree = \sqrt{\frac{1}{N}\sum_{j=1}^{N}
    \left(\frac{P_j^{\text{actual}}-P_j^{\text{ref}}}{P_j^{\text{ref}}}\right)^{\!2}}.
\end{equation}
$\Rthree=0$ indicates perfect dispatch fidelity; higher values indicate
greater deviation.  Unreachable generators contribute the maximum
relative deviation of~1.0.

These two metrics serve complementary diagnostic roles.  $\Rone$ answers
``how much load survives?''---an aggregate measure.  $\Rthree$ answers
``how accurately are surviving generators executing their dispatch?''---a
functional quality measure.  A system may retain high~$\Rone$ yet suffer
high~$\Rthree$, a form of hidden degradation that load-retention metrics
alone cannot detect.


%% ====================================================================
%%  SECTION 4 — EXPERIMENTAL ANALYSIS
%% ====================================================================
\section{Experimental Analysis}\label{sec:experiment}

%% ---- 4.1 Setup ------------------------------------------------------
\subsection{Experimental Setup}\label{sec:setup}

\paragraph{Test system.}
The physical layer is the IEEE~39-bus system (39~buses, 46~branches,
10~generators).  The cyber layer is a BA network with
$|V_C|=47$~nodes (39~non-CC + 8~CC), grown from a seed of
$m_0=4$ nodes with $m_e=2$ attachment edges.

\paragraph{Monte Carlo design.}
For each of the 1\,000 independent trials, a new BA network, a new CC
set (via MSIS), and a new $A_{pc}$ are generated.  The power layer is
shared across trials.

\paragraph{Parameter grid.}
Each trial is executed under all five delay scenarios at each of
11~tolerance levels ($\alpha\in\{0.0,0.1,\dots,1.0\}$), producing
$1{,}000\times5\times11=55{,}000$ independent cascade simulations.

The complete parameter set is given in Table~\ref{tbl:params}.

\begin{table}[t]
\caption{Simulation parameters.}\label{tbl:params}
\centering
\begin{tabular*}{\columnwidth}{@{\extracolsep{\fill}}lll@{}}
\toprule
Category & Parameter & Value \\
\midrule
\multirow{2}{*}{Power layer}
  & Test case & IEEE 39-bus \\
  & Buses / Branches / Gen. & 39 / 46 / 10 \\
\midrule
\multirow{3}{*}{Cyber layer}
  & Total nodes $|V_C|$ & 47 \\
  & Control centers $n_{\text{cc}}$ & 8 \\
  & BA seed $m_0$ / Attach.\ $m_e$ & 4 / 2 \\
\midrule
\multirow{2}{*}{Coupling}
  & Strategy & Betweenness-assortative \\
  & CC selection & MSIS ranking \\
\midrule
\multirow{3}{*}{Cascade}
  & Propagation prob.\ $p$ & 0.3 \\
  & Tolerance $\alpha$ & 0.0 : 0.1 : 1.0 \\
  & Attack mode & Betw.-targeted \\
\midrule
\multirow{5}{*}{\shortstack[l]{Communi-\\cation}}
  & Uplink packet $S_{\text{up}}$ & 8\,192 bits \\
  & Downlink packet $S_{\text{down}}$ & 2\,048 bits \\
  & Link rate $R$ & 10 Mbps \\
  & Prop.\ speed $V$ & $2\times10^5$ km/s \\
  & Link distance $d$ & 1 km \\
\midrule
\multirow{2}{*}{\shortstack[l]{Service\\delays}}
  & CC: tx / rx / fwd & 3 / 4 / 3 ms \\
  & Non-CC: tx / rx / fwd & 12 / 9 / 6 ms \\
\midrule
\multirow{4}{*}{\shortstack[l]{Delay-to-\\control}}
  & Meas.\ interface $\tau_{\text{PB}\to\text{nCC}}$ & 100 ms \\
  & Exec.\ interface $\tau_{\text{nCC}\to\text{PB}}$ & 120 ms \\
  & Meas.\ sensitivity $k_m$ & 0.80 s$^{-1}$ \\
  & Exec.\ sensitivity $k_e$ & 0.60 s$^{-1}$ \\
\midrule
Sampling & Trials per $(\alpha,\text{scen.})$ & 1\,000 \\
\bottomrule
\end{tabular*}
\end{table}

\paragraph{Evaluation protocol.}
For each cascade simulation, round-by-round snapshots record all failed
components, surviving topology, and the delay injection log (containing
$\eta_i$, $\tau_{m,i}$, $\tau_{e,i}$, reachability status for every
generator at every round).  Final-state $\Rone$ is computed from
surviving bus loads; final-state $\Rthree$ from the terminal-round delay
log.  Trials that terminate early contribute their terminal~$\Rone$ to
subsequent rounds via last-value-carried-forward (LVCF) padding.

%% ---- 4.2 Load Retention Under Delay ---------------------------------
\subsection{Impact of Delay Severity on Load Retention}\label{sec:R1result}

Fig.~\ref{fig:R1} presents the mean load retention ratio~$\Rone$ as a
function of tolerance factor~$\alpha$ across five delay scenarios,
averaged over 1\,000~trials.

\begin{figure}[t]
\centering
%% \includegraphics[width=\columnwidth]{fig_R1_vs_alpha}
\rule{\columnwidth}{4cm}  % placeholder
\caption{Mean load retention ratio~$\Rone$ versus tolerance
  factor~$\alpha$ across five delay scenarios.  The \nodelay{} curve
  provides an upper bound; the inter-scenario gap widens at
  moderate-to-high~$\alpha$.}\label{fig:R1}
\end{figure}

Three patterns emerge.  First, the \nodelay{} curve lies consistently
above all delay-present curves across the entire $\alpha$~range,
confirming that zero-delay control provides an upper bound on system
resilience.

Second, the gap between \nodelay{} and delay-present scenarios widens
dramatically as $\alpha$ increases through the intermediate range.  At
$\alpha=0.2$, all curves cluster near $\Rone\approx0.20$---the system
collapses regardless of delay.  By $\alpha=0.8$, the gap exceeds 0.25
in absolute terms.  This widening reveals the \emph{asymmetric
vulnerability}: delay damage is most consequential when the power system
has enough margin to survive without delay but not enough to absorb the
additional overloads that delay creates.

Third, the transition from high to low~$\Rone$ steepens under heavier
delay.  Delay effectively narrows the operating window within which the
system remains safe and shifts the critical transition point upward
along the tolerance axis.

%% ---- 4.3 Generator Dispatch Deviation --------------------------------
\subsection{Generator Output Fidelity Under Delay}\label{sec:R3result}

Fig.~\ref{fig:R3} presents $\Rthree$ across the same scenario--$\alpha$
grid.

\begin{figure}[t]
\centering
%% \includegraphics[width=\columnwidth]{fig_R3_vs_alpha}
\rule{\columnwidth}{4cm}  % placeholder
\caption{Generator dispatch deviation~$\Rthree$ versus~$\alpha$.  All
  five curves are perfectly stratified (no crossings), with \texttt{heavy}
  at the top and \nodelay{} at the bottom.}\label{fig:R3}
\end{figure}

The results are strikingly clean: all five curves are perfectly
stratified with no crossings, maintaining strict monotonic ordering from
\texttt{heavy} (worst fidelity) to \nodelay{} (best fidelity).

$\Rthree$ \emph{decreases} with increasing~$\alpha$---the opposite
direction from $\Rone$'s increase.  At higher~$\alpha$, fewer components
overload, fewer generators lose reachability, and surviving generators
retain higher~$\eta_i$ values.  The inter-scenario separation persists
even at high~$\alpha$ where $\Rone$ differences compress, revealing
hidden degradation: a system that retains most load yet exhibits
substantial control mismatch at the generator level.  Such systems
appear healthy by aggregate measures but are functionally stressed.

%% ---- 4.4 Distributional Analysis ------------------------------------
\subsection{Distributional Analysis of Cascade Outcomes}\label{sec:boxplot}

Fig.~\ref{fig:boxplot} presents box plots of the $\Rone$ distribution
at three representative tolerance levels ($\alpha=0.2,0.5,0.8$).

\begin{figure}[t]
\centering
%% \includegraphics[width=\columnwidth]{fig_R1_boxplot}
\rule{\columnwidth}{4cm}  % placeholder
\caption{Box plots of $\Rone$ distribution at $\alpha=0.2$, 0.5,
  and~0.8.  Delay amplifies distributional spread and introduces a
  catastrophic-collapse tail absent from the \nodelay{}
  case.}\label{fig:boxplot}
\end{figure}

At low~$\alpha$ ($=0.2$), all scenarios produce broadly similar
distributions---the system collapses regardless of delay.  At
moderate~$\alpha$ ($=0.5$), a striking divergence emerges: the \nodelay{}
distribution shifts upward while delay-present scenarios develop lower
whiskers extending toward $\Rone\approx0.05$, revealing a tail of
near-total collapse absent from \nodelay{}.  At high~$\alpha$ ($=0.8$),
\nodelay{} concentrates near $\Rone\approx1$ while delay-present
scenarios remain dispersed in the 0.40--0.60 range.

The mechanism underlying this variability amplification is the
interaction between random cyber topology and the delay feedback loop.
Under ideal communication, topological differences affect only which
nodes fail structurally.  Under heavy delay, topological differences
translate into different~$\eta_i$ profiles, different overload patterns,
and divergent cascade trajectories.  Small differences in cyber
connectivity, inconsequential under ideal communication, become decisive
under heavy delay.

%% ---- 4.5 Delay Penalty Heatmap --------------------------------------
\subsection{Temporal-Spatial Structure of the Delay Penalty}\label{sec:heatmap}

Define the round-level delay penalty:
\begin{equation}\label{eq:penalty}
  \Delta\Rone(r,\alpha)
  = \Rone^{\text{no\_delay}}(r,\alpha) - \Rone^{\text{heavy}}(r,\alpha).
\end{equation}

Fig.~\ref{fig:heatmap} presents $\Delta\Rone$ as a heatmap over the
(cascade round, $\alpha$) plane.

\begin{figure}[t]
\centering
%% \includegraphics[width=\columnwidth]{fig_heatmap}
\rule{\columnwidth}{4cm}  % placeholder
\caption{Delay penalty $\Delta\Rone$ heatmap.  The penalty ignites at
  Round~2, amplifies through Rounds~3--5, and saturates thereafter.  Low
  $\alpha$ remains dark (floor effect); the largest penalty concentrates
  at moderate-to-high~$\alpha$.}\label{fig:heatmap}
\end{figure}

The heatmap reveals a striking temporal-spatial structure:

\emph{Round~1 is delay-neutral.}  The leftmost column is uniformly dark
($\Delta\Rone\approx0$): Round~1 reflects the immediate structural
consequences of the initial attack, identical across delay scenarios.

\emph{The delay penalty ignites at Round~2.}  A sharp transition occurs
at $\alpha\geq0.4$: $\Delta\Rone$ leaps from near zero to substantial
values in a single round.  This is the moment the feedback loop first
fires.

\emph{Amplification through Rounds~3--5.}  The penalty intensifies as
each round of delay-induced overloads creates further structural damage,
producing compounding rather than linear accumulation.

\emph{Saturation after Round~5.}  Colors stabilize: 80--90\% of
avoidable delay damage is locked in by the fifth round.

\emph{Low $\alpha$ remains dark.}  When physical margins are so tight
that even the delay-free system collapses, there is no room for delay to
widen the gap.

\paragraph{Operational interpretation.}
An operator who could mitigate communication delay would achieve maximum
benefit by acting at Round~2---before the feedback loop amplifies the
initial penalty into compounding damage.  By Round~5, the intervention
window has largely closed.

%% ---- 4.6 Safety Classification --------------------------------------
\subsection{Safety-Level Classification}\label{sec:safety}

Based on~$\Rone$, each $(\alpha,\text{scenario})$ outcome is mapped to
an operational safety level (Table~\ref{tbl:safety}).

\begin{table}[t]
\caption{Safety-level classification.}\label{tbl:safety}
\centering
\begin{tabular}{@{}lcc@{}}
\toprule
Level & $\Rone$ range & Interpretation \\
\midrule
\textcolor{green!70!black}{Green}
  & $>85\%$ & Normal operation \\
\textcolor{yellow!80!black}{Yellow}
  & 70\%--85\% & Caution required \\
\textcolor{orange}{Orange}
  & 50\%--70\% & High risk \\
\textcolor{red}{Red}
  & $<50\%$ & System collapse \\
\bottomrule
\end{tabular}
\end{table}

The classification reveals that delay shifts the safety boundaries.  A
tolerance level yielding Green under \nodelay{} may yield Yellow or
Orange under \texttt{heavy} delay---not because the system is
structurally different, but because delay triggers additional cascade
rounds that the delay-free analysis would not predict.  Safety margins
computed without accounting for communication delay are therefore
potentially non-conservative.

%% ---- 4.7 Sensitivity Analysis (Planned) -----------------------------
\subsection{Sensitivity Analysis of Delay-Mitigation Actions}\label{sec:sensitivity}

Five candidate engineering actions are evaluated, each modifying a
single delay parameter under the \texttt{heavy} scenario while all
others remain at their heavy-delay values (Table~\ref{tbl:actions}).

\begin{table}[t]
\caption{Delay-mitigation actions.}\label{tbl:actions}
\centering
\begin{tabular*}{\columnwidth}{@{\extracolsep{\fill}}clp{2.8cm}@{}}
\toprule
Action & Description & Parameter change \\
\midrule
A1 & Bandwidth upgrade & $R$: 10\,$\to$\,100\,Mbps \\
A2 & Endpoint optimization & CC/non-CC tx/rx halved \\
A3 & Forwarding optimization & CC/non-CC fwd halved \\
A4 & Meas.\ interface speedup & 100\,$\to$\,20\,ms \\
A5 & Exec.\ interface speedup & 120\,$\to$\,30\,ms \\
\bottomrule
\end{tabular*}
\end{table}

For each action, the full Monte Carlo is repeated across all
$\alpha$~values.  The resulting $\Rone$ and $\Rthree$ distributions are
compared against both the unmodified \texttt{heavy} scenario and the
\nodelay{} benchmark to quantify each action's marginal benefit in
absolute terms ($\Delta\Rone$) and as a fraction of the maximum
recoverable penalty.

\emph{Note: Detailed numerical results will be incorporated once the
simulation completes.}

%% ---- 4.8 Timing-Action Interaction ----------------------------------
\subsection{Timing-Action Interaction Experiment}\label{sec:timing}

The heatmap analysis identified Round~2 as the critical intervention
window and $\alpha=0.5$--$0.6$ as the most delay-sensitive operating
regime.  A $2\times2$ factorial experiment crosses timing (correct:
Round~2; incorrect: Round~5) with action quality (best vs.\ worst
action from Section~\ref{sec:sensitivity}).  This tests whether the
timing-action interaction is additive or synergistic.

\emph{Note: Detailed numerical results will be incorporated once the
simulation completes.}


%% ====================================================================
%%  SECTION 5 — CONCLUSION
%% ====================================================================
\section{Conclusion}\label{sec:conclusion}

This paper constructs a delay-integrated cascading failure framework for
cyber-physical power systems in which communication delay is not an
after-the-fact observation but an endogenous causal variable embedded
within the cascade process.  At each cascade round, every surviving
generator's output is scaled by a delay-dependent control efficiency
factor~$\eta_i$ before the DC power flow that determines overloads.
This closes a feedback loop that prior models leave open: delay reduces
generation, altered power flows cause additional overloads, structural
failures degrade the communication topology, and the degraded topology
amplifies delay for subsequent rounds.

The framework models communication delay with full physical
grounding---decomposing end-to-end latency into serialization,
propagation, endpoint service, and forwarding components with
directional asymmetry between uplink measurement traffic and downlink
command traffic.  Control centers are placed via a multi-attribute
structural importance score (MSIS), and interdependent coupling uses a
betweenness-assortative strategy that pairs structurally critical power
buses with structurally critical communication nodes.

Monte Carlo simulation across 1\,000 independent coupling realizations,
11~tolerance levels, and 5~delay scenarios on the IEEE~39-bus system
yields four principal findings:

\begin{enumerate}
\item \textbf{Genuinely different cascade trajectories.}  Each delay
  scenario traces a unique path through the failure space---different
  overload patterns, failure sequences, and terminal states.  This
  qualitative divergence is fundamentally impossible in post-hoc delay
  models.

\item \textbf{Asymmetrically concentrated delay penalty.}  The
  $\Delta\Rone$ heatmap reveals that delay damage is largest at
  moderate-to-high tolerance margins and late cascade rounds---precisely
  where delay-free analysis predicts safety.  Delay-free models are
  most misleading where operators are most likely to rely on them.

\item \textbf{Amplified stochastic variability.}  Under heavy delay,
  the $\Rone$ distribution widens substantially, introducing a
  catastrophic-collapse tail absent from the ideal-communication case.

\item \textbf{Hidden functional degradation.}  $\Rthree$ maintains
  clean stratification across delay scenarios even at high~$\alpha$
  where $\Rone$ differences compress, exposing a functional quality
  deficit invisible to aggregate metrics.
\end{enumerate}

The practical implication is direct: robustness assessment and
safety-margin setting for CPPS must incorporate communication delay as a
causal variable within the cascade model.  Defense investments should
prioritize maintaining communication reachability---through path
redundancy, CC diversity, and connectivity-preserving
reinforcement---over minimizing per-hop latency, because complete loss
of controllability ($\eta=0$) produces qualitatively more severe damage
than degraded controllability ($0<\eta<1$).  The round-level analysis
further suggests that delay-mitigation interventions are most effective
early in the cascade (Round~2) before the feedback loop compounds
initial damage into irreversible collapse.

Future work will extend this framework in three directions: (i)~a
nonlinear or threshold-based delay-to-control mapping that captures
stability loss beyond critical delay values; (ii)~generalization to
larger test systems (IEEE~118-bus) and alternative communication
topologies; and (iii)~delay-aware defense design that explicitly targets
the reachability-dominated fragmentation threshold.


%% ====================================================================
%%  References
%% ====================================================================
\bibliographystyle{cas-model2-names}
\bibliography{paper2_refs}

\end{document}
